\chapter{Conclusion}\label{ch:conclusion}
\graphicspath{{../Figures/80-conclusion/}{figs/80-conclusion/}}

This dissertation moved the biomimetic humanoid robot project from an unvalidated design concept to a quantitatively characterized actuator set and experimentally tested knee-torque models, and it defined the path from those experimental foundations to neuromechanical control.

At the actuator level, the analysis in Chapter~\ref{ch:muscles} produced novel equations for calculating isometric force in Festo \qtylist{10;20}{\mm}-diameter BPAs. The work elucidated the relationship between maximum BPA force at \qty{620}{\kPa} and resting length, $F_{620}(l_0)$, a previously unreported characteristic of these actuators that the manufacturer's tool predicts with the opposite trend, and it produced a nondimensional force surface, $F^*(\epsilon^*,P^*)$, with only three diameter-dependent coefficients. Taken together, these fits predict isometric BPA force across resting lengths from a small number of measurements, correct the large over-predictions of the manufacturer's tool at short resting lengths, and outperform published single-length models. In a design workflow, these equations make it possible to select actuator diameter and resting length and to evaluate whether a proposed muscle routing can meet force and range-of-motion requirements, before committing to a prototype.

At the joint level (Chapter~\ref{ch:torque}), the BPA model was combined with transmission corrections for effective length, series compliance, and wrapping. Identification on the pinned knee improved torque prediction, and the shared correction terms also improved predictions on the biomimetic four-bar knee. The corrected model supports route redesign: the optimized two-BPA flexor and cam-extensor configurations are predicted to meet their human-scaled torque targets across the design range. The isometric tests demonstrate that the redesigned configurations meet or exceed the human reference torques over most of the tested range of motion.

The Sensory Afferent Database in Chapter~\ref{ch:database} organizes physiological evidence by properties and experimental context, making it easier to select and assess candidate feedback pathways. Chapter~\ref{ch:sns} addresses the complementary problem of testing those pathways in coupled neural and mechanical models. I developed the Li--Deng walker implementations, Gait2392-based neuromechanical models, SNS-Simscape neural blocks and actuator integrations, and staged tuning and experiment procedures. Selected demonstrations establish implementation agreement for a contact-driven rhythm core, corrective feedback in a loaded-arm model, and feedback-dependent stepping in a fixed supported W2L configuration. These results show that the tools can distinguish circuit behavior, sensory contribution, and mechanical response under specified conditions.

The demonstrated contributions support controlled neuromechanical experiments with explicit neural pathways and actuator descriptions. Extending these results to unsupported whole-body behavior and hardware requires additional balance, dynamic-actuation, and transfer studies (Chapter~\ref{ch:overall}). The present findings concern the measured actuator and joint behavior and the specified simulation tests.

The dissertation provides experimentally grounded actuator and joint models, a curated afferent resource, and functioning neuromechanical modeling and evaluation tools. Together, they make actuator selection, torque prediction, sensory-pathway selection, and controlled simulation tests more explicit and reproducible. These are usable contributions to biomimetic robot design and to studying how neural control interacts with mechanical structure.


